\section{Nine natural unknowns with a smaller degree}
\label{sec:nine-improved}

The fewest-variable construction already present in package 18~\cite{p18} has nine
nonnegative unknowns. We keep its explicit quartic-to-Jones specialization,
replace its final relation-combining step, and use bounds adapted to the
intermediate expressions. The resulting fixed polynomial has degree
\[
 \boxed{93264\cdot5^{33}+288
 =10\,857\,358\,574\,867\,248\,535\,156\,250\,288}
\]
and a supplied division-free circuit with $20\,547$ binary operations.
The previous degree was $161844\cdot5^{33}+484$, at $20\,531$ operations.
Thus the degree decreases by approximately $42.3741\%$, for sixteen
additional operations. The number of unknowns remains nine. This is an
improvement of the supplied fixed-instance representation; it is not a new
nine-variable representation theorem or a global minimality assertion.

\subsection{The inherited quartic, specified completely}

To make every coefficient determinate, we use the literal
\texttt{quartic32.upstream.json} of package 18~\cite{p18}. In this subsection only,
its positive coordinates are
\begin{align*}
 &(J,F,\alpha,z,f,h,i,o,j,s,w,\tau,\eta,\zeta,y,\gamma,\\[-2pt]
 &\hspace{35mm}N,E,P,c,a,X,D,C_2,G,L,F_2,V,g,I,B_{\mathrm{aux}},\Delta).
\end{align*}
The four fixed ports are the entire Wang943 tuple
$\kappa=(\beta,K_0,M_C,M_F)$; its fourth port is the shifted mask.
The affine abbreviations are
\[
 q=\beta J+1,\qquad k=\eta+\zeta,\qquad H=4a+3,
 \qquad C_w=q-F-\alpha.
\]
The following 23 expressions are the residuals $r_j$:
\begin{center}
\small
\begin{longtable}{@{}r@{\quad}l@{}}
\toprule $j$&$r_j$\\\midrule
\endfirsthead
\toprule $j$&$r_j$\\\midrule
\endhead
1&$L-q^2+1$\\
2&$N-q(L+1)$\\
3&$X-wq$\\
4&$a-E-sN$\\
5&$E-X(a-E)$\\
6&$c-k(a-E)-\eta$\\
7&$P-E(c-\eta)-k$\\
8&$\tau^2-P(P-k)-1$\\
9&$D-X-ac-\gamma H$\\
10&$\Delta-a^2-H$\\
11&$C_2-c^2$\\
12&$D^2-\Delta C_2-1$\\
13&$k-jc+V-1-hE$\\
14&$G-(q-1)(q-F)-\alpha$\\
15&$jc-V-LG-(M_C+qM_F)J$\\
16&$(K_0+w)C_w+q-F-z(q-1)-1$\\
17&$F_2-f^2$\\
18&$V-of+c$\\
19&$g-y^2+V^2$\\
20&$I-iC_2$\\
21&$B_{\mathrm{aux}}-I^2$\\
22&$B_{\mathrm{aux}}-\Delta(F_2-1)$\\
23&$y^2-B_{\mathrm{aux}}g-1$\\
\bottomrule
\end{longtable}
\end{center}
All residuals have degree at most two in the 32 supplied coordinates.
Define, with a dummy input $x$ which does not occur,
\[
 \mathcal P_{32}(x,\mathbf n)=\sum_{j=1}^{23}r_j(\mathbf n+\mathbf1)^2,
 \qquad \mathbf n\in\mathbb N^{32}.
\]
It is a quartic: the coefficient of $n_\tau^4$ is one. It satisfies the
normalization $\mathcal P_{32}(x,\mathbf0)\ge1$, because residual 5 equals
one at the all-ones positive tuple.

The semantic dependency is the inherited ordinary-strong quartic theorem:
these positive residual equations are equivalent in existence to the
inputless native/Pell equations. That theorem is a four-port interface
result~\cite{p18,arithmetic}. Package 19~\cite{p19} proves that its \emph{entire} Wang943 tuple implements
the same interface, including its canonical positive completion. Hence it
can replace the original T4 tuple in all 23 residuals simultaneously.
No conclusion is drawn from mask congruences alone, and no individual port
is transferred in isolation. The ordinary compiler theorem and K215
calibration then give
\[
 \exists\mathbf n\in\mathbb N^{32}\;\mathcal P_{32}(1,\mathbf n)=0
 \quad\Longleftrightarrow\quad \neg\operatorname{Con}(\mathrm{ZFC}).
\]
This remains an ordinary mathematical proof dependency, not a new
proof-assistant endpoint.

The literal row 16 uses the old transport quotient. Some inherited notes
write its shifted form $(K_0+w)C_w=F+z_{\rm new}(q-1)$; the conventions are
related by $z=z_{\rm new}+1$. Canonical witnesses have $z\ge2$, as proved
in the inherited positive-completion argument. The present coefficient
recipe uses precisely the residual table above, without silently changing
that convention. Inherited relative-path metadata inside the immutable
source JSON is historical; the new complete binding is the Wang943 tuple
specified in this report and its binding file.

\subsection{Bounded definitions of every Jones coefficient}

Put $\nu=32$ and $L=5^{33}$. For
$\alpha=(\alpha_0,\ldots,\alpha_{32})\in\mathbb N^{33}$ with
$|\alpha|\le4$, define
\begin{align*}
 p_\alpha&=[x^{\alpha_0}n_1^{\alpha_1}\cdots n_{32}^{\alpha_{32}}]
                 \mathcal P_{32},\\
 c_\alpha&=(4-|\alpha|)!\prod_{j=0}^{32}\alpha_j!\;p_\alpha,
 &e(\alpha)&=\sum_{j=0}^{32}\alpha_j5^j.
\end{align*}
The factorial of the unused degree is part of the construction. There are
$\binom{37}{4}=66\,045$ eligible multi-indices. The delivered data record
346 nonzero coefficient polynomials in the four fixed ports.

After expanding each affine abbreviation, write
\[
 r_j(\mathbf v)=\sum_{a,b}d_{j,a,b}\mathbf v^a\kappa^b,
 \qquad d_{j,a,b}\in\mathbb Z.
\]
The following is a finite, entirely specified integer recipe:
\[
 \rho_j=\sum_{a,b}|d_{j,a,b}|\,2^{|a|}\kappa^b,
 \quad s_\star=2+24\sum_{j=1}^{23}\rho_j^2,
 \quad z_\star=2^{s_\star}.
\]
All four ports are positive. The coefficient $\ell^1$ norm after replacing
$\mathbf v$ by $\mathbf n+\mathbf1$ is at most $\rho_j$; this norm is
submultiplicative, and every factorial multiplier is at most 24. Therefore
\[
 \sum_{|\alpha|\le4}|c_\alpha|\le s_\star-2,
 \qquad |c_\alpha|+1<z_\star.
\]
Finally put
\begin{align*}
 e_\star(T)&=\sum_{|\alpha|\le4}(z_\star+c_\alpha)T^{L-e(\alpha)},\\
 u_\star&=\sum_{j=1}^{32}(2z_\star)^{5^j},
 &y_\star&=e_\star(2z_\star).
\end{align*}
Since $e(\alpha)\le4\cdot5^{32}<L$, every exponent is a fixed natural
number. All digits are positive. These are the exact index constants in
the constructive Jones specialization; $u_\star$ enters its proof
interface but is not a leaf of the final circuit. There is no search for
an unknown recursively enumerable index. The index integers are not
materialized: their exact bounded recipes and modular evaluators are the
delivered representation.

\subsection{All nine-variable polynomial expressions}

From here onward the symbols are local to the Jones construction. The
nine supplied unknowns are exactly
\[
 (\varepsilon,g_0,h_0,s_0,w_0,\phi,i_0,j_0,n)\in\mathbb N^9.
\]
Set $g=g_0+1$, $h=h_0+1$, $s=s_0+1$, $w=w_0+1$,
$i=i_0+1$, and $j=j_0+1$. Define
\begin{align*}
 b&=y_\star+1+\varepsilon,&\mathcal B&=b^5,&q&=\mathcal B^L,\\
 \ell&=\sum_{r=1}^{32}\mathcal B^{5^r},&
 \lambda&=\sum_{r=0}^{4L-1}\mathcal B^r,&e&=e_\star(\mathcal B),\\
 \theta&=\mathcal B-2z_\star,&N&=q^{16},&
 G_{\rm cap}&=32b\mathcal B^{5^{32}}.
\end{align*}
The packed quantity is
\begin{align*}
R={}&\left[g+eq^3+\ell q^5+
 \left(2(e-z_\star\lambda)(1+\mathcal B+g)^4
       +\lambda\mathcal B+\lambda\mathcal Bq^4\right)q^7\right](N^2-N)\\
 &+\left[q^3-b\ell+\ell+\theta\lambda q^3+
                   (\mathcal B-2)q^8\right](N^2-1).
\end{align*}
Continue with
\begin{align*}
 U&=N^2w,&Y&=N^2s,&M&=RY,\\
 P&=2M^2U,&K&=R+1+h(P-1),&A&=M(U+1),\\
 B_1&=2R+1,&W_1&=bw,&C&=B_1+W_1+\phi,\\
 J&=4A-5,&Q&=3W_1C-2(W_1^2-1),&D&=(A^2-1)C^2+1,\\
 E&=iJDC^2,&F&=(A^2-1)E^2+1,&G_1&=A+F(CD+1-A),\\
 H&=B_1+jC,&I&=(G_1^2-1)H^2+1.&
\end{align*}
The two radicands, signed divisor, dividend, and margin are
\begin{align*}
 a_0&=(P^2-1)K^2+1,&a_1&=DFI,\\
 u&=JF,&v&=QF+J(H-C),\\
 m&=K^2G_{\rm cap}-4(C-KY)^2G_{\rm cap}-8gK^2.
\end{align*}
All displayed quantities are polynomial abbreviations, not additional
unknowns. In particular $R$ is not assumed positive.

\subsection{Unconditional bounds before imposing any equation}

\begin{lemma}[The packed expression never vanishes]
\label{lem:nine-R-nonzero}
On every nonnegative assignment of the eight outer coordinates,
$R\ne0$, $|A|\ge2$, $P\ge2$, and $K\ge1$.
\end{lemma}
\begin{proof}
The positive index recipe gives $b\ge2$. Reducing the displayed formula
for $R$ modulo $q$ gives
\[
 R\equiv(b-1)\ell\pmod q.
\]
Indeed $q\mid N$, and all terms in its second bracket except
$(1-b)\ell$ are divisible by $q$. Every exponent occurring in $\ell$
belongs to $\{0,\ldots,L-1\}$, so
\[
 0<(b-1)\ell\le(\mathcal B-1)\ell
       \le(\mathcal B-1)\sum_{r=0}^{L-1}\mathcal B^r=q-1<q.
\]
Thus $R\ne0$. Put $\sigma=Y(U+1)$. Since $U,Y\ge1$, we have
$\sigma\ge2$ and $A=R\sigma$, whence $|A|\ge2$.
Moreover $P=2R^2Y^2U\ge2R^2\ge2$, and
\[
 K=R+1+h(P-1)\ge R+P\ge2R^2+R\ge1.
\]
None of these conclusions uses a square condition, divisibility condition,
or a sign inferred from $m$.
\end{proof}

Set
\begin{align*}
 r&=W_1+\phi,&\sigma&=Y(U+1),&T_C&=(j+2)C^2+r^2,\\
 M_D&=\sigma(R^2+C^2),&M_F&=AE+1,\\
 M_I&=FD T_C+2FH^2+(F-1)A^2+1,\\
 V_0&=PK+1,&V_1&=M_D M_F M_I+1.
\end{align*}
The circuit uses the equal, slightly cheaper expression
$M_I=F(DT_C+2H^2+A^2)-A^2+1$.

\begin{lemma}[Bounds adapted to the three Pell radicands]
\label{lem:nine-structured-bounds}
For every assignment under consideration,
\[
 V_0,V_1\ge2,\qquad
 \sqrt{|a_0|}\le V_0-1,\qquad
 \sqrt{|a_1|}\le V_1-1.
\]
Also $J\ne0$, $D,F\ge1$, and therefore $u=JF\ne0$.
\end{lemma}
\begin{proof}
The preceding lemma gives $|A|\ge2$, so $J=4A-5\ne0$ and
$AJ=A(4A-5)>0$. Consequently $D\ge1$,
$AE=iAJDC^2\ge0$, and $F\ge1$.
Since $K\ge1$,
\[
 1\le a_0=(PK)^2-K^2+1\le(PK)^2,
\]
which proves the $V_0$ assertion.

For $D$, substitute $A=R\sigma$ and calculate
\[
 M_D^2-D=\sigma^2(R^4+R^2C^2+C^4)+C^2-1
       \ge\sigma^2R^4-1\ge0.
\]
Also $M_D\ge1$. Since $AE\ge0$,
\[
 1\le F=(AE)^2-E^2+1\le(AE+1)^2=M_F^2,
 \qquad M_F\ge1.
\]
To bound $I$ without an unstated sign hypothesis, use
\[
 H=(j+1)C-r,\qquad G_1=FDC+F-A(F-1).
\]
The elementary inequalities $|rC|\le r^2+C^2$,
$|H|\le H^2$ for integer $H$, and $|AH|\le A^2+H^2$ give
\begin{align*}
 |CH|&\le(j+1)C^2+r|C|\le T_C,\\
 |G_1H|&\le FD T_C+FH^2+(F-1)(A^2+H^2).
\end{align*}
Thus $M_I\ge|G_1H|+|H|+1$, and
\[
 \sqrt{|I|}\le\sqrt{G_1^2H^2+H^2+1}
             \le|G_1H|+|H|+1\le M_I.
\]
All three majorants are at least one, including when $C=0$ or $H=0$.
Multiplying the bounds for $D,F,|I|$ proves
$\sqrt{|DFI|}\le M_D M_F M_I=V_1-1$.
\end{proof}

\subsection{A signed-divisor relation-combining lemma}

The second improvement avoids squaring both the divisor and dividend.
The useful condition is a bound on the rational quotient, which is already
available in this application.

\begin{lemma}[Combining two squares, signed divisibility, and positivity]
\label{lem:signed-divisor-combiner}
Let $a_0,a_1,u,v,m,M,V_0,V_1$ be integers satisfying
\[
 u\ne0,\quad V_0,V_1\ge2,\quad
 \sqrt{|a_t|}\le V_t-1\ (t=0,1),\quad
 M\ge1,\quad |v/u|\le M-1.
\]
For $n\in\mathbb N$, put
\[
 \Omega=u(2m-1),\qquad
 \Theta=u(n+M)+v-\Omega(2M+V_0V_1)
\]
and
\[
 \mathcal C=
 \bigl(\Theta^2-\Omega^2(a_0+V_0^2a_1)\bigr)^2
                  -4\Omega^4V_0^2a_0a_1.
\]
Then
\[
 \exists n\in\mathbb N\;\mathcal C=0
 \quad\Longleftrightarrow\quad
 a_0,a_1\text{ are integer squares},\quad u\mid v,\quad m>0.
\]
\end{lemma}
\begin{proof}
Because $u\ne0$ and $2m-1\ne0$, also $\Omega\ne0$.
Choose complex square roots $r_t^2=a_t$. Multiplication gives
\[
 \mathcal C=
 \prod_{\epsilon,\delta\in\{-1,1\}}
       \bigl(\Theta-\Omega(\epsilon r_0+\delta V_0r_1)\bigr).
\]
If this vanishes, absorb the signs of a vanishing factor into the roots.
Then $s=\Theta/\Omega=r_0+V_0r_1\in\mathbb Q$. If $s\ne0$, the identity
$2sr_0=s^2+a_0-V_0^2a_1$ proves both roots rational, hence integral
because their squares are integers. If $s=0$, the bounds imply
$|r_1|=|r_0|/V_0\le(V_0-1)/V_0<1$. The integer $a_1$ therefore has
absolute value less than one, so $a_1=a_0=0$. This also gives integral
roots. No independence of quadratic fields is assumed.

Set $Z=2M+V_0V_1+r_0+V_0r_1$. The root bounds give
\[
 Z\ge2M+V_0V_1-(V_0-1)-V_0(V_1-1)=2M+1>0.
\]
Dividing the vanishing linear factor by the possibly negative $u$ yields
\[
 n+M+v/u=(2m-1)Z.
\]
Its left side is at least $n+1>0$. Hence $m>0$; its right side is an
integer, so $v/u$ is an integer and $u\mid v$.
Conversely, under the three stated conditions choose integer roots and put
\[
 n=(2m-1)Z-M-v/u.
\]
This is an integer and
$n\ge2M+1-M-|v/u|\ge2$. Its linear factor vanishes, proving the
converse. This proves the lemma for both signs of $u$.
\end{proof}

In the Jones application choose the polynomial
\[
 M_{\rm div}=Q^2+(H-C)^2+1.
\]
The identity
\[
 \frac vu=\frac QJ+\frac{H-C}{F}
\]
and $J\ne0$, $F\ge1$ give
$|v/u|\le|Q|+|H-C|\le M_{\rm div}-1$.
No division occurs in the final polynomial; the rational identity is used
only in its proof.

\subsection{The resulting equation and consistency equivalence}

Define
\begin{align*}
 \Omega&=JF(2m-1),\\
 \Theta&=JF(n+M_{\rm div})+QF+J(H-C)
               -\Omega(2M_{\rm div}+V_0V_1).
\end{align*}
\[
 \boxed{\mathcal F_9^{\rm new}
   =\bigl(\Theta^2-\Omega^2(a_0+V_0^2a_1)\bigr)^2
                   -4\Omega^4V_0^2a_0a_1.}
\]
Every symbol on the right was defined by finite polynomial formulas. The
requested equation in this section is $\mathcal F_9^{\rm new}=0$ over
$\mathbb N^9$.

\begin{theorem}[The improved nine-variable specialization]
With the exact coefficients fixed above, ZFC proves
\[
 \bigl(\nexists\mathbf x\in\mathbb N^9:
                \mathcal F_9^{\rm new}(\mathbf x)=0\bigr)
 \quad\Longleftrightarrow\quad\operatorname{Con}(\mathrm{ZFC}).
\]
\end{theorem}
\begin{proof}
The new unconditional bounds and signed-divisor lemma show that,
for each eight-coordinate outer assignment, existence of the ninth
coordinate is equivalent to the same four Jones ratio conditions as
in package 18: $a_0,a_1$ are squares, $JF\mid QF+J(H-C)$, and $m>0$.
The inherited constructive Jones lemmas
\texttt{Nine.ratio\_complete} and \texttt{Nine.ratio\_sound}, composed
with their packing reduction~\cite{jones,ninemain,p18}, identify existence of those eight
coordinates with a nonnegative zero of $\mathcal P_{32}(1,\cdot)$.
Our explicit normalized quartic and index recipes discharge their
hypotheses exactly as in package 18. The ordinary quartic/compiler and
K215 calibration then identify such a zero with a finite ZFC refutation.
Negating the resulting existence equivalence proves the theorem.
The change is in the finalizer and its bounds, not in the inherited
ratio theorem. The full new endpoint has not been kernel checked.
\end{proof}

\subsection{Exact degree and the operation ledger}

All fixed ports and Jones index constants have degree zero. For a direct
certificate, specialize every supplied unknown to $t$, keeping all fixed
leaves symbolic. The following table gives exact degrees and leading
coefficients on that line; the gate-by-gate calculation has no unresolved
leading cancellation.
\begin{center}
\small
\begin{tabular}{@{}lrr@{}}
\toprule Expression&Degree&Leading coefficient\\\midrule
$R$&$235L$&$1$\\
$P$&$950L+3$&$2$\\
$K$&$950L+4$&$2$\\
$A$&$555L+2$&$1$\\
$a_0$&$3800L+14$&$2^4$\\
$a_1$&$24640L+68$&$2^{46}$\\
$JF$&$6875L+20$&$2^{14}$\\
$m$&$2221L+11$&$-2^9$\\
$M_D$&$790L+2$&$5$\\
$M_F$&$3160L+9$&$2^6$\\
$M_I$&$8370L+23$&$2^{16}$\\
$V_0$&$1900L+7$&$2^2$\\
$V_1$&$12320L+34$&$5\cdot2^{22}$\\
$\Omega$&$9096L+31$&$-2^{24}$\\
$\Theta$&$23316L+72$&$5\cdot2^{48}$\\
\bottomrule
\end{tabular}
\end{center}
Here the two leading terms inside the outer square have equal degree,
so one must check their coefficients instead of adding degree bounds
blindly. Their difference has leading coefficient
\[
 (25-4)2^{96}=21\cdot2^{96}.
\]
The last term $4\Omega^4V_0^2a_0a_1$ has lower degree. Therefore the
final leading coefficient is $441\cdot2^{192}\ne0$, independently of
all fixed constants, proving the exact degree $93264L+288$.

The delivered live circuit has 668 additions, 23 subtractions, and
$19\,856$ multiplications, totaling $20\,547$ operations. It expands
every fixed power into multiplication gates. It has exactly nine
variable leaves and 348 nonliteral fixed leaves: $z_\star,y_\star$ and
the 346 nonzero $c_\alpha$. These are fixed integers with bounded
recipes. Their bit lengths are not included in the operation count.

The independent verification reads the actual four Wang943 integers and
checks their hashes, reconstructs the normalized quartic and all 346
coefficient records, and compares the literal DAG to direct finite sums
modulo three primes. It also verifies the exact leading coefficient,
6,912 structured-majorant assignments including negative $R$, $C=0$ and
$H=0$, and 134,400 bounded signed-finalizer evaluations, with 468 explicitly
constructed completeness witnesses. These audits support implementation
and transcription; they do not replace the universal proofs above or
decide whether the polynomial has a solution.

